\scanpage{202}%
\mspage{76}%
culièrement \unsure{crucials}, le cas d'une intersection $C\cap W$ !

Le résultat principal, pour lequel j'ai fait les \unsure{gammes} du § précédent, est le suivant :

\uline{Th. 5.1} Soit $M$ une catégorie satisfaisant les conditions dites :

% slanted
\margin{Démonstration \unsure{canulée|canular}, Énoncé peut-être \uuline{faux} !!} a) Stabilité par petites $\underrightarrow{\lim}$ (pas néc. filtr.

b) $M$ muni d'une filtration cardinale $(\mathrm{Filt}^{\unsure{\pi|R}})_{\pi\geq\pi_0}$,

c) les objets de $M$ sont accessibles

d) les $\mathrm{Filt}^\pi$ sont stables par $\underrightarrow{\lim}\!,$ indexées par des cat. $I$ telles que $\operatorname{Card}\mathrm{Fl}(I)\leq\pi$, et si $Z\in\mathrm{Filt}^\pi$, tout quotient $Z'$ de $Z$ (tel que $Z\longrightarrow Z'$ ne se factorise pas par un sous-objet $\unsure{Z'_0|Z'}$ de $Z'$ $\ne Z'$) est dans $\mathrm{Filt}^\pi$.

e) les limites inductives filtrantes dans $M$ transforment monos en monos.

\margin{$\big|$} Soit \deleted{\unsure{donc}} $C$ une partie L-accessible de $\mathrm{Fl}(M)=\mathrm{Ob}\,\underline{\mathrm{Fl}}\,(M)$, i.e. \inserted{\unsure{(5.17)|(5.11)}} une partie stable par $\underrightarrow{\lim}$ filtrantes, et accessible. \deleted{\unsure{Alors}} \inserted{Supposons de plus $C$, formé de monomorphismes\({}^{\ast}\), et \unsure{en} contenant les \uline{Iso}.}
% slanted
\margin{$\ast$ finalement, on a \unsure{eu besoin de|besoin de} cette hypothèse} Alors \deleted{soit} il existe une partie petite $C_0\subset C$, telle que $C\subset\widetilde C_0$. Plus précisément, si $\pi$ est un \deleted{\unsure{plus}} cardinal tel que $C$ soit $\pi$-accessible, et si on pose $C^\pi=C\cap\mathrm{Fl}^\pi M$, \deleted{de sorte qu'il} \inserted{alors on a}
% archive sheet number: 201
